\section{Results}\label{sec:results}

\subsection{The same design-level fault is not the same physical fault}\label{sec:res-severity}

For the same tap pair the SCIG's fault current is about 10\,\% lower than the PMSG's (medians 3.3 against 3.7\,A for inter-turn and 7.7 against 8.5\,A for inter-winding pairs), but the difference that matters for detection is relative to what the stator sensors see: the fault current is 1.45 times the pre-fault stator current (RMS to RMS, median over pairs) in the PMSG and 0.63 times in the SCIG for inter-turn faults, and 3.3 against 1.5 times for inter-winding faults, because the SCIG's stator current carries the magnetising component (5.3 against 2.6\,A RMS before the fault). The WFSG trials at the matching impedance give 1.20 (inter-turn, 11.6--12.1\,\%) and 1.03 (inter-winding, 2.3--10.6\,\%), and 0.65 at the lower impedance used for its small inter-turn extents. Being resistor-limited (Section~\ref{sec:extent}), these currents scale with volts per turn and speed; a bolted fault would reach several times rated current and---the first-order difference between the topologies---could be de-excited in the SCIG and the WFSG but not in the PMSG. A fault specified as ``$e$\,\% of a winding'' therefore arrives at the sensors of the three alternatives as three different physical events; Section~\ref{sec:res-mde} uses both axes.

\subsection{Where the signature appears, and why the SDR differs}\label{sec:res-where}

The null of the SDR differs by up to a factor of 35 between the PMSG and the SCIG (Table~\ref{tab:decomp}, columns $q_{95}$), and not only in sensed quantities: the null of the duty-cycle negative sequence, a controller output that involves no sensor, is nine times larger on the PMSG, so the healthy non-stationarity is a property of the machine--drive combination, not of the transducers alone. The table also decomposes the SCIG/PMSG ratio of median SDR: for inter-turn faults the PMSG's relative changes are the \emph{larger} ones on every feature (SCIG/PMSG ratio of median $|\delta|$ 0.37--0.71), and the SCIG's higher SDR is entirely a quieter null (ratio of null quantiles 2.6--35); for inter-winding faults the numerators are comparable (0.3--1.4) and the denominators again decide.

\begin{table}[t]
\scriptsize
\setlength{\tabcolsep}{3pt}
\caption{SCIG/PMSG ratio of median SDR decomposed into numerator and denominator, five-cycle windows}\label{tab:decomp}
\begin{tabular}{@{}llccccccc@{}}
\toprule
Class & Feature & $|\delta|$ P & $|\delta|$ S & $q_{95}$ P & $q_{95}$ S & Num. & Den. & SDR ratio \\
\midrule
inter-turn & $V_2/V_1$ & 0.257 & 0.136 & 0.086 & 0.033 & 0.53 & 2.64 & 1.40 \\
inter-turn & $I_2/I_1$ & 0.238 & 0.152 & 0.162 & 0.034 & 0.64 & 4.82 & 3.08 \\
inter-turn & $I_d$\,2$f_e$ & 0.291 & 0.107 & 0.563 & 0.016 & 0.37 & 35.09 & 12.95 \\
inter-turn & PI$_d$\,2$f_e$ & 0.289 & 0.107 & 0.557 & 0.016 & 0.37 & 34.71 & 12.91 \\
inter-turn & $V_q$\,2$f_e$ & 0.222 & 0.155 & 0.165 & 0.060 & 0.70 & 2.73 & 1.91 \\
inter-turn & $V_q^{\mathrm{cmd}}$\,2$f_e$ & 0.197 & 0.139 & 0.162 & 0.039 & 0.71 & 4.17 & 2.94 \\
inter-turn & $D_2/D_1$ & 0.251 & 0.111 & 0.090 & 0.010 & 0.44 & 8.96 & 3.98 \\
inter-winding & $V_2/V_1$ & 5.318 & 1.569 & 0.086 & 0.033 & 0.30 & 2.64 & 0.78 \\
inter-winding & $I_2/I_1$ & 1.833 & 1.869 & 0.162 & 0.034 & 1.02 & 4.82 & 4.91 \\
inter-winding & $I_d$\,2$f_e$ & 2.335 & 1.188 & 0.563 & 0.016 & 0.51 & 35.09 & 17.85 \\
inter-winding & PI$_d$\,2$f_e$ & 2.344 & 1.188 & 0.557 & 0.016 & 0.51 & 34.71 & 17.60 \\
inter-winding & $V_q$\,2$f_e$ & 3.688 & 1.495 & 0.165 & 0.060 & 0.41 & 2.73 & 1.11 \\
inter-winding & $V_q^{\mathrm{cmd}}$\,2$f_e$ & 0.845 & 1.184 & 0.162 & 0.039 & 1.40 & 4.17 & 5.84 \\
inter-winding & $D_2/D_1$ & 2.942 & 0.854 & 0.090 & 0.010 & 0.29 & 8.96 & 2.60 \\
\botrule
\end{tabular}
\footnotetext{Num.: SCIG/PMSG ratio of the median relative change under fault $|\delta|$; Den.: ratio of null quantiles $q_{95}^{\mathrm{PMSG}}/q_{95}^{\mathrm{SCIG}}$; P: PMSG; S: SCIG.}
\end{table}


Fig.~\ref{fig:features} shows the median SDR of every feature for the three alternatives (the WFSG pooled over its fault impedances; features to the right of the dashed line at $\mathrm{SDR}=1$ are detectable in more than half of the trials). In the PMSG the terminal-voltage group carries the strongest signature in 67\,\% of inter-turn and 88\,\% of inter-winding trials, through $V_2/V_1$ (median SDR 3.0 and 62 at five-cycle windows; 76 and 94\,\% of trials detectable), and the stator-current negative sequence $I_2/I_1$ is weak for inter-turn faults (median SDR 1.5, 56\,\%). In the SCIG the controller group is strongest in 56 and 67\,\% of trials, through the 2$f_e$ ripple of the $d$-axis current and its PI output (6.7 and 74; 79 and 100\,\%), the stator-current group in 28--33\,\% and the terminal-voltage group in 13\,\% and none. The two machines rank the 30 common features with a Spearman correlation of 0.70 (inter-turn) and 0.89 (inter-winding): the same features are informative in both, in a different order. In the WFSG the strongest inter-turn features are $V_2/V_1$, the $q$-axis current ripple and the commanded $q$-axis voltage, and for inter-winding faults the $q$-axis ripple, the commanded voltage, the duty cycles and the third current harmonic; its field-current 2$f_e$ ripple, an observable the other two machines lack, detected 86\,\% of the inter-turn trials at the matching impedance and all inter-winding trials (Table~\ref{tab:suites3}). The physical quantity behind the relocation is the incremental negative-sequence impedance $\Delta V_2/\Delta I_2$ seen at the terminals under inter-winding faults: 15 [11, 18]\,$\Omega$ in the PMSG, 7.0 [6.6, 7.5]\,$\Omega$ in the SCIG and 11 [10, 13]\,$\Omega$ in the WFSG (medians and interquartile ranges over trials, three-cycle windows), i.e.\ the PMSG drive converts a given negative-sequence current disturbance into twice the terminal voltage that the SCIG drive does; for inter-turn faults the PMSG's negative-sequence current increments are of mixed sign with a median indistinguishable from zero, which is why current-side features see little there. Both machines run the same field-oriented controller but differ in angle source (encoder against flux observer), negative-sequence impedance and feature baselines (healthy medians of $I_2/I_1$, $V_2/V_1$ and $D_2/D_1$: 0.0055, 0.0009 and 0.0014 on the PMSG against 0.0079, 0.0020 and 0.0036 on the SCIG, and 0.0148, 0.0073 and 0.0039 on the WFSG), so the relocation belongs to the machine--drive combination, not to the control structure alone.

\begin{figure}[t]
\centering
\includegraphics[width=\textwidth]{G_sdr_by_feature_3alt}
\caption{Median SDR of every feature (dot) with interquartile range (bar) for the three alternatives, three-cycle windows}\label{fig:features}
\end{figure}

\subsection{Minimum detectable extent}\label{sec:res-mde}

Table~\ref{tab:mde} gives the minimum detectable extent with its bootstrap interval, the detectable share with its Wilson interval and the median SDR with its tap-pair bootstrap interval for four representative features; Fig.~\ref{fig:severity} plots the SDR against the design-level extent (left; dots are trials, large markers the median per extent level, their shape the fault impedance) and against the physical severity axis (right; RMS fault current relative to the RMS pre-fault stator current, all trials). Through the terminal-voltage negative sequence the PMSG detects inter-turn faults from 3.8\,\% [2.8, 3.8] and the SCIG from 7.4\,\% [7.4, 7.4]; through the stator-current negative sequence both need 7.4\,\%, and through the $d$-axis ripple the PMSG needs 12.1\,\% and the SCIG 7.4\,\% [3.8, 7.4]. The smallest tested inter-turn level (2.7--2.8\,\%) is not reliably detectable on either machine by any single feature. For inter-winding faults the SCIG detects from the smallest tested extent (2.3\,\%) on every feature except the $q$-axis voltage ripple, while the PMSG needs 2.3\,\% [2.3, 14.3] on $V_2/V_1$ but 9.7\,\% on $I_2/I_1$ and 21.7\,\% on the $d$-axis ripple. The SDR rises with extent in every machine (Spearman 0.51--0.95), and the MDE intervals are wide because the tested extents form three clusters (about 3, 7.5 and 12\,\% for inter-turn faults), so the SCIG/PMSG ratio of median SDR is the better basis for comparison: 3.1 [1.3, 5.6] for $I_2/I_1$ and 12 [2.6, 30] for the $d$-axis ripple on inter-turn faults, but 1.1 [0.4, 4.6] for $V_2/V_1$ and 1.5 [0.5, 3.2] for the $q$-axis voltage ripple. The two machines differ where the signature relocates and are indistinguishable where it does not. On the physical axis (Fig.~\ref{fig:severity}, right) the picture is the same and independent of the tap-pair design: at equal fault-to-stator current ratio the SCIG's current-side features lie above the WFSG's, which lie above the PMSG's, while on $V_2/V_1$ the three machines nearly coincide. For the WFSG, whose small inter-turn extents were tested only at the lower impedance, the design-level MDE at the matching impedance is bounded by the smallest extent tested there (11.6\,\%, detectable in 78--98\,\% of trials depending on the suite); its 2.7--7.7\,\% inter-turn trials at 1\,$\Omega$ sit at SDR $\approx1$ on $V_2/V_1$ and $I_2/I_1$, and all its inter-winding extents from 2.3\,\% are detectable.

\begin{table}[t]
\scriptsize
\setlength{\tabcolsep}{2pt}
\caption{Minimum detectable extent, detectable share and median SDR with 95\,\% intervals for four features, PMSG and SCIG, five-cycle windows, $\alpha=0.95$}\label{tab:mde}
\begin{tabular}{@{}lllcccc@{}}
\toprule
Class & Feature & M. & MDE [CI] & DS [CI] & Median SDR [CI] & S/P ratio [CI] \\
\midrule
IT & $V_2/V_1$ & P & 3.8 [2.8, 3.8] & 76 [67, 83] & 2.99 [1.21, 12.5] & \\
 &  & S & 7.4 [7.4, 7.4] & 73 [64, 81] & 4.17 [0.949, 15.5] & 1.1 [0.378, 4.56] \\
IT & $I_2/I_1$ & P & 7.4 [7.4, 12.1] & 56 [47, 65] & 1.47 [0.601, 2.36] & \\
 &  & S & 7.4 [7.4, 7.4] & 75 [66, 82] & 4.52 [1.11, 11.7] & 3.07 [1.34, 5.63] \\
IT & PI$_d$\,2$f_e$ & P & 12.1 [11.6, 12.1] & 31 [23, 40] & 0.518 [0.389, 1.11] & \\
 &  & S & 7.4 [3.8, 7.4] & 78 [69, 85] & 6.69 [1.3, 18.3] & 11.8 [2.59, 28.8] \\
IT & $V_q$\,2$f_e$ & P & 7.4 [7.4, 11.6] & 57 [48, 66] & 1.35 [0.791, 4.58] & \\
 &  & S & 7.4 [7.4, 7.7] & 60 [51, 69] & 2.58 [0.516, 5.33] & 1.49 [0.464, 3.21] \\
IW & $V_2/V_1$ & P & 2.3 [2.3, 14.3] & 94 [88, 97] & 61.9 [9.24, 181] & \\
 &  & S & 2.3 [2.3, 2.3] & 100 [97, 100] & 48.2 [15.3, 188] & 0.974 [0.696, 1.68] \\
IW & $I_2/I_1$ & P & 9.7 [9.5, 21.7] & 81 [73, 88] & 11.3 [1.98, 64.3] & \\
 &  & S & 2.3 [2.3, 2.3] & 100 [97, 100] & 55.6 [12.5, 117] & 5.49 [1.75, 13.1] \\
IW & PI$_d$\,2$f_e$ & P & 21.7 [9.7, 21.7] & 72 [63, 80] & 4.21 [1.18, 13] & \\
 &  & S & 2.3 [2.3, 2.3] & 100 [97, 100] & 74 [22.9, 245] & 18.2 [12.3, 27.2] \\
IW & $V_q$\,2$f_e$ & P & 9.7 [2.3, 14.3] & 90 [83, 94] & 22.4 [3.24, 62] & \\
 &  & S & 9.7 [2.3, 9.7] & 92 [85, 96] & 24.8 [7.04, 78.6] & 1.13 [0.831, 1.9] \\
\botrule
\end{tabular}
\footnotetext{MDE: \% of the winding branch between the taps, limit-of-detection convention, bootstrap over operating points within each extent level; DS: \%, Wilson interval; median SDR and SCIG/PMSG ratio: bootstrap over tap pairs, the same pairs for both machines. IT: inter-turn; IW: inter-winding; P: PMSG; S: SCIG. `--': the largest tested extent is not detectable; smallest tested extents 2.7\,\% (IT) and 2.3\,\% (IW).}
\end{table}


\begin{figure}[t]
\centering
\includegraphics[width=\textwidth]{G_severity_combined}
\caption{SDR of $V_2/V_1$ (\textbf{a}--\textbf{d}) and $I_2/I_1$ (\textbf{e}--\textbf{h}) against the design-level extent (left) and the physical severity (right), three alternatives, three-cycle windows}\label{fig:severity}
\end{figure}

\subsection{Operating envelope}\label{sec:res-op}

Speed dominates the operating-point dependence: in the PMSG the median SDR of the best feature rises from 5--9 at 1200\,rpm to 16--21 at 1800\,rpm and the detectable share from 79--83 to 92--96\,\%, in the SCIG from 16--18 to 44--51 and from 92--96 to 96--100\,\%; torque changes little. Two cautions apply: the dependence is consistent with a resistor-limited fault current that scales with speed, so it is a property of the emulation as much as of the machines; and the pooled null is a mixture over operating points (the share of trials whose healthy change exceeds the pooled 95th percentile is 0--16\,\% by speed for the key features), so the 5\,\% false-alarm reading holds on average over the envelope. Both make the lowest monitored speed a parameter of the requirement.

\subsection{Observation suites and the design decision}\label{sec:res-suites}

Table~\ref{tab:suites} gives the suite results for the PMSG and SCIG at five-cycle windows, Table~\ref{tab:suites3} (Appendix~\ref{app:appendix}) for the three alternatives on the common feature set at three cycles, and Table~\ref{tab:decision} the decision table of step~5. For the PMSG, current transducers alone resolve inter-turn faults from 7.4\,\% and inter-winding faults from 9.7\,\%; three voltage transducers take these to 3.8 and 2.3\,\%, and neither controller access nor a torque transducer improves the fixed-feature result further. The drive-internal suite reaches 7.4\,\% for inter-turn faults through the duty-cycle negative sequence, a weaker version of the signature the measured terminal voltage carries. For the SCIG the cheapest suites already reach what any suite reaches---7.4\,\% for inter-turn faults through $I_2/I_1$ or, with no added sensor, the $d$-axis PI output, and 2.3\,\% for inter-winding faults---and additional hardware raises the detectable share rather than the reach. For the WFSG every suite detects all inter-winding faults from 2.3\,\% and 78--98\,\% of the inter-turn trials at the matching impedance, and the field current alone detects 86 and 100\,\%. The nested selection of the fixed feature changes no share by more than five points (Table~\ref{tab:suites3}); the oracle shares come with false-alarm rates of 44--78\,\% on the healthy trials for the large suites and are upper bounds only. A shaft torque transducer adds nothing for winding faults on the two machines on which it was available (the torque features' healthy relative change exceeds 100\,\% on the SCIG and their median SDR is below 0.3 on the PMSG).

\begin{table}[t]
\tabsize
\setlength{\tabcolsep}{3pt}
\caption{Observation suites versus what they detect, PMSG and SCIG, five-cycle windows: best fixed feature, its MDE (\%) and detectable share (\%)}\label{tab:suites}
\begin{tabular}{@{}lllllll@{}}
\toprule
Suite & PMSG feature & MDE & DS [oracle] & SCIG feature & MDE & DS [oracle] \\
\midrule
\multicolumn{6}{@{}l}{\emph{Inter-turn faults}} \\
drive-internal & $D_2/D_1$ & 7.4 & 64 [87] & PI$_d$\,2$f_e$ & 7.4 & 78 [99] \\
3 CT & $I_2/I_1$ & 7.4 & 56 [77] & $I_2/I_1$ & 7.4 & 75 [87] \\
3 CT + 3 VT & $V_2/V_1$ & 3.8 & 76 [88] & $I_2/I_1$ & 7.4 & 75 [90] \\
3 CT + 3 VT + drive & $V_2/V_1$ & 3.8 & 76 [95] & PI$_d$\,2$f_e$ & 7.4 & 78 [99] \\
+ torque transducer & $V_2/V_1$ & 3.8 & 76 [95] & PI$_d$\,2$f_e$ & 7.4 & 78 [99] \\
\multicolumn{6}{@{}l}{\emph{Inter-winding faults}} \\
drive-internal & $D_2/D_1$ & 14.3 & 90 [100] & PI$_d$\,2$f_e$ & 2.3 & 100 [100] \\
3 CT & $I_2/I_1$ & 9.7 & 81 [94] & $I_2/I_1$ & 2.3 & 100 [100] \\
3 CT + 3 VT & $V_2/V_1$ & 2.3 & 94 [97] & $I_2/I_1$ & 2.3 & 100 [100] \\
3 CT + 3 VT + drive & $V_2/V_1$ & 2.3 & 94 [100] & PI$_d$\,2$f_e$ & 2.3 & 100 [100] \\
+ torque transducer & $V_2/V_1$ & 2.3 & 94 [100] & PI$_d$\,2$f_e$ & 2.3 & 100 [100] \\
\botrule
\end{tabular}
\footnotetext{Oracle: per-trial best feature, its detectable share in brackets (false-alarm rates in Online Resource~1, Table~S1); features failing the reliability screen are excluded for that machine.}
\end{table}


The decision table turns this into the selection of step~5. At $\alpha=0.95$ no suite on any alternative meets an inter-turn requirement of 3 or 5\,\% with 95\,\% bootstrap confidence; the PMSG's voltage suite is the only candidate below 5\,\% and its verdict is uncertain, a consequence of the three-cycle windows and the coarse extent grid. At 8\,\% the drive-internal suite meets the requirement on the PMSG and the SCIG, i.e.\ no added sensor is needed for either; the WFSG's inter-turn verdicts are bounded by its smallest tested extent. For inter-winding faults the SCIG and the WFSG meet every requirement from 3\,\% with the drive-internal suite, whereas the PMSG needs current transducers at 12\,\% and is uncertain below (its five-cycle result, Table~\ref{tab:suites}, reaches 2.3\,\% with voltage transducers). Relaxing $\alpha$ to 0.90 admits the PMSG's voltage suite at 5\,\% for inter-turn faults; tightening it to 0.99 removes the PMSG's current-only option at 12\,\% for inter-winding faults.

\begin{table}[t]
\tabsize
\setlength{\tabcolsep}{3pt}
\caption{Decision table: cheapest suite whose MDE meets $e_{\mathrm{req}}$ with probability $\ge0.95$, per alternative, fault class and $\alpha$}\label{tab:decision}
\begin{tabular}{@{}llcccccc@{}}
\toprule
\multirow{2}{*}{$\alpha$} & \multirow{2}{*}{$e_{\mathrm{req}}$ (\%)} & \multicolumn{2}{c}{PMSG} & \multicolumn{2}{c}{SCIG} & \multicolumn{2}{c}{WFSG} \\
\cmidrule(lr){3-4}\cmidrule(lr){5-6}\cmidrule(l){7-8}
 & & IT & IW & IT & IW & IT & IW \\
\midrule
0.90 & 3 & (3CT+3VT)? & (drive)? & none & drive & none & drive \\
0.90 & 5 & 3CT+3VT & (drive)? & none & drive & none & drive \\
0.90 & 8 & drive & (drive)? & drive & drive & none & drive \\
0.90 & 12 & drive & 3CT & drive & drive & drive & drive \\
0.95 & 3 & (3CT+3VT)? & (drive)? & none & drive & none & drive \\
0.95 & 5 & (drive)? & (drive)? & none & drive & none & drive \\
0.95 & 8 & drive & (drive)? & drive & drive & none & drive \\
0.95 & 12 & drive & 3CT & drive & drive & drive & drive \\
0.99 & 3 & none & none & none & drive & none & drive \\
0.99 & 5 & (3CT+3VT)? & none & none & drive & none & drive \\
0.99 & 8 & drive & none & drive & drive & none & drive \\
0.99 & 12 & drive & (drive)? & drive & drive & drive & drive \\
\botrule
\end{tabular}
\footnotetext{Three-cycle windows; probabilities from 1000 bootstrap resamples. Suites ordered drive $<$ 3CT $<$ 3CT+3VT $<$ 3CT+3VT+drv; exc: WFSG field current; IT: inter-turn; IW: inter-winding. `(suite)?': no suite meets with that confidence and the cheapest suite with an uncertain verdict is shown; `none': every suite fails. WFSG inter-turn verdicts are bounded by its smallest tested extent at the matching impedance (11.6\,\%).}
\end{table}


\subsection{Transfer across alternatives}\label{sec:res-transfer}

Table~\ref{tab:transfer} reports step~6 for the PMSG--SCIG pair and Table~\ref{tab:transfer3} (Appendix~\ref{app:appendix}) for the three alternatives; the figures quoted below are AUCs on the out-of-reference windows, and the values on all windows agree with them to within 0.02, so the windows used for normalisation do not drive the result. In physical units a gradient-boosted detector reaches an area under the receiver operating characteristic curve (AUC) of 0.87--0.91 within a machine on unseen tap pairs but 0.78 (PMSG$\to$SCIG) and 0.55 (SCIG$\to$PMSG) across; with its features scaled on the target machine's healthy trials, a commissioning baseline, 0.82 and 0.75, and a logistic detector under that rule falls below chance in one direction (0.45); with each feature referenced to the first half of the trial's own reference interval, 0.91 and 0.84, against 0.91--0.94 within a machine, and a logistic detector transfers at 0.89 and 0.87. The threshold that gives 1\,\% false positives on all negatives fires on 0--3\,\% of the healthy trials' fault-time windows, the only negatives that contain the contactor operation. Between the three alternatives, self-referenced detectors transfer in all six directions with AUC 0.80--0.98 on the out-of-reference windows, and a detector trained on two machines detects the faults of the third with 0.83 (PMSG), 0.88 (SCIG) and 0.97 (WFSG; optimistic, as that bench has no healthy trials). Restricted to each suite's features, every suite transfers at 0.82--0.95 when self-referenced and at 0.62--0.90 under the commissioning baseline.

\begin{table}[t]
\scriptsize
\caption{Transfer of a window-level detector between PMSG and SCIG: AUC on all windows / AUC on the out-of-reference windows}\label{tab:transfer}
\begin{tabular}{@{}llcccc@{}}
\toprule
Referencing & Model & P$\to$P & S$\to$S & P$\to$S & S$\to$P \\
\midrule
physical units & logistic & 0.78/0.78 & 0.79/0.80 & 0.71/0.71 & 0.59/0.59 \\
physical units & GBDT & 0.88/0.87 & 0.90/0.91 & 0.78/0.78 & 0.55/0.55 \\
$z$ on target's healthy trials & logistic & 0.78/0.78 & 0.79/0.80 & 0.71/0.71 & 0.44/0.45 \\
$z$ on target's healthy trials & GBDT & 0.88/0.87 & 0.90/0.91 & 0.82/0.82 & 0.75/0.75 \\
$|z|$ vs own reference & logistic & 0.90/0.89 & 0.94/0.93 & 0.90/0.89 & 0.88/0.87 \\
$|z|$ vs own reference & GBDT & 0.92/0.91 & 0.95/0.94 & 0.92/0.91 & 0.86/0.84 \\
\botrule
\end{tabular}
\footnotetext{Out-of-reference windows: second half of the reference interval, recovery windows and fault-time windows of the healthy trials. P$\to$P, S$\to$S: five-fold cross-validation grouped by tap pair (healthy trials by speed); P$\to$S, S$\to$P: trained on all trials of one machine, tested on all trials of the other.}
\end{table}


\subsection{Robustness and the two nulls}\label{sec:res-robust}

Table~\ref{tab:sens} (Appendix~\ref{app:appendix}) shows how the tuning choices move the design quantities, with the screen applied at every setting. The inter-turn MDEs of all five key features are the same at three, five and eight cycles on both machines and the detectable shares move by at most six points; for inter-winding faults one MDE moves (the PMSG's $V_2/V_1$, 9.7\,\% at three cycles against 2.3\,\% at five and eight). Tightening the null quantile from 0.95 to 0.99 lowers the shares by 3--9 points and moves two PMSG inter-winding MDEs up one level; relaxing it to 0.90 raises the shares by up to 9 points and moves one down. No ordering of the two machines reverses, but the PMSG's inter-winding suite conclusion depends on $\alpha$ and on the window length (Table~\ref{tab:decision}). At $\alpha=0.95$ the screen removes up to three weak features per machine and window length (Appendix~\ref{app:computation}); none of the five key features is removed at any setting, and the three-cycle result of $V_2/V_1$ on the PMSG (two false alarms in nine healthy trials) passes the binomial rule. An absolute-change variant of the SDR, which removes the healthy level of the feature from numerator and null, ranks the 30 features with a Spearman correlation of 0.92--0.98 against the relative variant and picks the same three top features in every machine and class. The bootstrap intervals (Table~\ref{tab:mde}) put the median SDRs within a factor of two to four and the shares within $\pm$8--10 points.

Table~\ref{tab:between} (Appendix~\ref{app:appendix}) replaces the reference immediately before the fault by the healthy trial recorded at the same operating point at another time, a commissioning baseline: the null of the four key features widens by a factor of 1.6 to 4.9, the median SDR of the inter-turn features falls to about one (PMSG $V_2/V_1$ from 3.0 to 1.1; SCIG $I_2/I_1$ from 4.5 to 1.1) and their detectable shares to 20--59\,\%, while inter-winding faults remain detectable in 67--94\,\% of trials. The MDEs reported here are therefore statements about a short-horizon change detector; a monitor that compares against an old baseline needs the wider null.
